\chapter[Inference and Reproducible Workflows]{Inference, Modelling, and Reproducible Line Workflows}
\label{ch:inference}

A measured line flux is the endpoint of a reduction model and the input to a
physical model.  Treating it as a context-free number loses calibration
covariance, nondetections, selection, response information, and the assumptions
used to separate line from continuum.  A reproducible workflow instead carries
a generative model from latent source spectrum to detector samples, exposes the
approximation used at each compression step, and stores enough provenance to
repeat the result.

\section{Begin with the data-generating process}

Let $\boldsymbol{s}(\boldsymbol{\theta})$ be a rest-frame source spectrum for
physical parameters $\boldsymbol{\theta}$, $A$ an aperture or spatial mixing
operator, $R$ the spectral response, and $B\boldsymbol{\phi}$ a background model.
A general expectation in detector space is
\begin{equation}
 \boldsymbol{\mu}(\boldsymbol{\theta},\boldsymbol{\phi})
 =R\,A\,\boldsymbol{s}(\boldsymbol{\theta})+B\boldsymbol{\phi}.
 \label{eq:forward-model}
\end{equation}
For photon counts $n_i$, the elementary likelihood is
\begin{equation}
 \ln \mathcal{L}_{\rm P}=\sum_i
 \left[n_i\ln\mu_i-\mu_i-\ln\Gamma(n_i+1)\right].
 \label{eq:poisson-likelihood}
\end{equation}
Background counts measured in another region should be modeled jointly with
their area and exposure scaling, not subtracted and assigned an ad hoc Gaussian
error.  This is the logic behind count-space statistics used in X-ray
spectroscopy \citep{Cash1979,Arnaud1996}.

For approximately Gaussian extracted spectra $\boldsymbol{y}$ with covariance
$C$, the corresponding log likelihood is
\begin{equation}
 \ln\mathcal{L}_{\rm G}=-\frac{1}{2}
 \left[(\boldsymbol{y}-\boldsymbol{m})^{\mathsf T}C^{-1}
 (\boldsymbol{y}-\boldsymbol{m})+\ln|C|+N\ln(2\pi)\right].
 \label{eq:gaussian-covariance}
\end{equation}
Resampling, sky subtraction, continuum fitting, common calibration, PSF
matching, and shared denominators create covariance.  Ignoring off-diagonal
terms can narrow intervals and exaggerate the information content of a line
profile.  If full $C$ is too large, preserve a sparse covariance, a correlation
kernel, or the linear operator that generated the spectrum.  Merely inflating
diagonal errors does not reproduce parameter-dependent correlations.

Equation~\ref{eq:forward-model} is also a reminder not to ``deconvolve first''
unless deconvolution is itself part of the likelihood.  A model profile should
normally be convolved with the wavelength-dependent line-spread function and
integrated over detector pixels.  For an integral-field cube, the spatial PSF,
slice geometry, differential atmospheric refraction, and source velocity field
can couple spatial and spectral structure.  Beam smearing can then mimic a
broad component or flatten a velocity gradient.  Forward-modeling the native
samples avoids asking a noisy inverse operation to create information.

\subsection{Fit the continuum and line together when they interact}

A local line model might include a polynomial or physically motivated continuum,
one or more response-convolved profiles, and nuisance terms for sky residuals.
Continuum uncertainty is especially important for weak equivalent widths and
broad lines.  Estimating a continuum, subtracting it, and then fitting the line
as if the continuum were known discards covariance.  Simultaneous fitting or
posterior propagation preserves it.

Line parameters need constraints that represent physics without manufacturing
detections.  A doublet may share velocity and width and have a ratio bounded by
atomic physics; blended species need not share kinematics.  Positive-only flux
priors change the null distribution at zero.  Consequently, the improvement in
$\chi^2$ from adding a non-negative line does not necessarily follow the usual
asymptotic likelihood-ratio law.  Calibrate searches with posterior predictive
checks or simulations that reproduce the continuum fit and look-elsewhere scan
\citep{Protassov2002}.

\subsection{Identifiability, priors, and parameterization}

A parameter is identifiable only to the extent that distinct values predict
distinguishable data after nuisance parameters and response are considered.
For an optically thin single line, column density, emitting area, abundance,
and ion fraction may appear only as a product.  No sampler can separate them
without additional information.  Report identifiable combinations, add
well-motivated complementary data, or acknowledge prior dominance.

Priors should be specified in the parameterization where the scientific belief
is intended.  Uniform priors in flux, log flux, equivalent width, and line
amplitude divided by noise are not equivalent.  A bounded ``uninformative''
prior can set an implicit physical scale through its endpoints.  Transforming
from positive width $\sigma$ to $\log\sigma$ changes the density by a Jacobian.
Prior predictive simulation - drawing parameters from the prior and generating
synthetic spectra - is an efficient way to expose absurd widths, fluxes,
doublet ratios, or continua before seeing the data.

Parameterization also affects computation.  Centered hierarchical models can
mix poorly when individual lines are weak, whereas a non-centered form can be
more stable.  Highly covariant amplitude and width may benefit from parameters
closer to integrated flux and log width.  Convergence diagnostics should be
reported for the actual posterior draws: multiple chains, effective sample
size, rank-based diagnostics, trace inspection, and sensitivity to initialization.
A maximum-likelihood covariance matrix is a local curvature approximation and
can miss boundaries, tails, multiple modes, and discrete alternatives.

\subsection{Model comparison without a significance shortcut}

Deciding between one and two kinematic components is a scientific model-
comparison problem, not simply a request for the model with lower residual sum
of squares.  Information criteria approximate predictive tradeoffs under
specific regularity assumptions; marginal likelihoods depend on the full prior
volume; cross-validation evaluates predictive performance for a declared unit
of held-out data.  Mixture components and lines at a boundary violate common
asymptotic assumptions.  The preferred tool therefore depends on the claim.

When the question is whether a broad outflow component is required, simulate
spectra from the best adequate narrow-component model, repeat the same search
over centroid and width, and compare the chosen detection statistic with the
simulated distribution.  When the question is prediction for new spectra,
leave out independent objects or exposures rather than correlated pixels.
When physical alternatives are all approximations, present model-averaged or
model-stratified results instead of converting a small score difference into
certainty about ontology.

\begin{tcolorbox}[title={Do not confuse a fit with a measurement},colback=red!3,colframe=red!45!black]
A converged optimizer establishes neither identifiability nor calibration.
Inspect alternative starting points, parameter correlations, residual structure,
prior sensitivity, and simulation-based recovery.  Quote the observable and its
likelihood or posterior, not only the best-fit parameter vector.
\end{tcolorbox}

\section{Ratios, calibration terms, and uncertainty budgets}

For a ratio $r=F_1/F_2$, first-order propagation gives
\begin{equation}
 \mathrm{Var}(\ln r)=
 \frac{\mathrm{Var}(F_1)}{F_1^2}+
 \frac{\mathrm{Var}(F_2)}{F_2^2}-
 2\frac{\mathrm{Cov}(F_1,F_2)}{F_1F_2}.
 \label{eq:ratio-covariance}
\end{equation}
The covariance term is often beneficial when both lines share flux calibration,
aperture correction, or continuum placement.  It can also be negative after
deblending.  Treating catalogue columns as independent can therefore either
overstate or understate ratio uncertainty.

A useful hierarchical calibration model for observation $j$ is
\begin{equation}
 F_{ij}^{\rm obs}=g_j\,a_{ij}\,F_i^{\rm true}+\epsilon_{ij},
 \qquad \ln g_j\sim\mathcal{N}(0,\sigma_{g,j}^2),
 \label{eq:calibration-nuisance}
\end{equation}
where $g_j$ is a shared multiplicative calibration and $a_{ij}$ represents
line-dependent attenuation or aperture effects.  Marginalizing these nuisance
terms propagates correlations naturally.  Keep statistical repeatability,
absolute calibration, model discrepancy, atomic data, distance, lensing,
attenuation, and aperture terms separate in the budget; they average down in
different ways.

Avoid constructing a diagnostic analysis entirely from logarithmic ratios when
one or more denominator lines are weak.  The ratio likelihood becomes skewed
and heavy-tailed, and a catalogue cut on both numerator and denominator creates
a complicated truncated sample.  A safer population model predicts the vector
of line fluxes, applies attenuation and calibration, and compares that vector
with all measured flux estimators and their covariance.  Ratios can still be
plotted as summaries.  The model, however, retains the information in the
overall normalization and admits negative noise realizations without undefined
logarithms.

\begin{table}[htbp]
\centering
\caption{A minimum uncertainty ledger for a derived line quantity.}
\label{tab:uncertainty-ledger}
\begin{tabularx}{\textwidth}{@{}p{0.22\textwidth}p{0.29\textwidth}X@{}}
\toprule
Class & Example & Recommended representation \\
\midrule
Counting/noise & source, sky, read noise & detector-space likelihood or covariance \\
Reduction & flat field, sky basis, extraction & nuisance parameters or reduction ensembles \\
Calibration & throughput, beam, standard star & shared multiplicative/additive terms \\
Model & continuum, profiles, phase mixture & competing models and posterior predictive checks \\
Physics data & collision strengths, abundances & versioned grid and sensitivity reruns \\
Population & distance, lensing, selection & hierarchical latent variables and selection function \\
\bottomrule
\end{tabularx}
\end{table}

\section{Nondetections and censored information}

An upper limit is not a flux measurement with value zero.  If an analysis
retains a Gaussian estimator $\hat F$ and uncertainty $\sigma_F$, use its full
likelihood even when $\hat F/\sigma_F$ fails a reporting threshold.  If only the
event $\hat F<F_{\rm lim}$ is available, its likelihood under a model predicting
$F(\boldsymbol{\theta})$ is
\begin{equation}
 p(\hat F<F_{\rm lim}\mid\boldsymbol{\theta})=
 \Phi\!\left(\frac{F_{\rm lim}-F(\boldsymbol{\theta})}{\sigma_F}\right),
 \label{eq:censored-gaussian}
\end{equation}
where $\Phi$ is the normal cumulative distribution.  This censored likelihood
is less informative than the original estimator.  Publishing signed fluxes and
uncertainties, alongside quality flags, generally supports better population
inference than publishing limits alone.

Classical survival-analysis methods are valuable when their censoring assumptions
hold \citep{FeigelsonNelson1985}.  Astronomical censoring is often informative:
the limit varies with wavelength, background, source size, and target selection,
all of which can correlate with the physical quantity.  Then the measurement
and selection process must be modeled jointly.  Stacking is not a universal
escape.  A mean stack estimates a weighted population functional set by the
normalization, weights, and missing data; it does not recover the median object
or the ratio of mean latent luminosities without assumptions.

\section{Selection effects and population likelihoods}

Suppose objects have latent parameters $\boldsymbol{\theta}$ drawn from
$p(\boldsymbol{\theta}\mid\boldsymbol{\Lambda})$ and enter the sample with
probability $P_{\rm det}(\boldsymbol{\theta})$.  Conditioning on a detected
sample introduces the efficiency
\begin{equation}
 \alpha(\boldsymbol{\Lambda})=
 \int P_{\rm det}(\boldsymbol{\theta})
 p(\boldsymbol{\theta}\mid\boldsymbol{\Lambda})\,d\boldsymbol{\theta}.
 \label{eq:selection-efficiency}
\end{equation}
For fixed detected sample size, a schematic population likelihood includes one
factor of $\alpha^{-1}$ per object,
\begin{equation}
 \mathcal{L}(\boldsymbol{\Lambda})\propto
 \prod_{k=1}^{N}\frac{
 \int p(d_k\mid\boldsymbol{\theta}_k)
 p(\boldsymbol{\theta}_k\mid\boldsymbol{\Lambda})\,d\boldsymbol{\theta}_k}
 {\alpha(\boldsymbol{\Lambda})}.
 \label{eq:population-selection}
\end{equation}
The exact form changes if the total event rate is modeled, but the principle is
stable: selection normalizes the population model \citep{MandelFarrGair2019}.

Injection and recovery is the most direct way to estimate a complicated
$P_{\rm det}$.  Inject sources spanning flux, width, wavelength, continuum,
size, and environment into realistic data; run the exact detection and quality
pipeline; and model recovery probability in that multi-dimensional space.
Injecting only isolated Gaussian lines into white noise misses sky residuals,
blends, continuum errors, and human vetting.  Avoid applying a one-dimensional
``90 percent flux completeness'' curve to a population whose width or size
distribution evolves.

\section{Physical models: useful approximations and model discrepancy}

\subsection{Photoionized nebulae and slabs}

Photoionization codes solve ionization, thermal balance, and radiative transfer
through a specified geometry and incident spectrum.  Cloudy is widely used for
nebulae, AGN, and diffuse gas \citep{Ferland2017}.  A grid must record code
version, atomic-data set, incident spectral energy distribution, gas density or
pressure law, abundances and depletion, dust prescription, stopping criterion,
geometry, covering factor, and whether emission is integrated or emergent from
a specified face.  ``Metallicity'' alone is not a complete model coordinate.

Interpolation across a grid should be tested against direct runs, especially
near ionization fronts or thermal instabilities.  Grid priors are not innocuous:
a uniform prior in $U$ differs from one uniform in $\log U$, and a rectangular
grid can assign weight to chemically inconsistent abundance combinations.
Model discrepancy should be included when observed line ratios cannot be matched
within measurement errors; otherwise the posterior may become extremely precise
about the least-wrong point in an inadequate model.

Large physical grids are often replaced by interpolators or machine-learning
emulators.  Validate the emulator on withheld grid points using errors in the
space actually fit - often log luminosity or line ratio - and propagate emulator
uncertainty.  Dense validation near sharp ionization fronts and grid edges is
more valuable than a small global root-mean-square error.  Never extrapolate
silently: mark proposals outside the training hull, expand the physical grid,
or assign a discrepancy wide enough to express ignorance.  Archive the training
grid, preprocessing, emulator weights, and code version alongside posterior
products.

\subsection{Collisionally ionized and nonequilibrium plasma}

Hot-plasma emissivities depend on electron temperature, ion fractions,
abundances, and atomic rates.  Databases such as CHIANTI provide versioned
atomic data and software for optically thin plasma calculations
\citep{Dere1997,DelZanna2021}.  Collisional ionization equilibrium assumes that
ionization and recombination have had time to balance at one temperature.
Cooling flows, shocks, interfaces, and rapidly heated gas can be in
nonequilibrium; ion populations then remember thermal history.  Resonant
scattering, charge exchange, non-Maxwellian electrons, or optical depth can
break an otherwise plausible one-temperature fit.

The observed X-ray count model must be folded through the response and include
foreground absorption.  A multi-temperature plasma fit by one temperature can
bias abundances.  Comparing multiple plasma codes or atomic-data releases is a
practical model-systematics test, not an optional software preference.

\subsection{From line emission to mass or rate}

Every conversion should be written as a conditional model.  Examples include
$M_{\rm mol}=\alpha_{\rm CO}L'_{\rm CO}$, an ionized-gas mass proportional to
$L_{\rm recomb}/(n_e\alpha_{\rm eff})$, and a star-formation calibration based
on ionizing-photon production.  The conversion coefficient depends on physical
state, composition, stellar population, IMF, leakage, and geometry.  Report the
direct observable alongside the derived value so later work can replace the
calibration without reverse engineering it \citep{KennicuttEvans2012}.

\section{Validation: can the inference recover known truth?}

Posterior predictive checks draw replicated data from the fitted model and
compare features relevant to the claim: residual autocorrelation, line wings,
doublet consistency, nondetection fraction, or ratio distribution.  A good
global $\chi^2$ can coexist with a systematic failure in the diagnostic of
interest.  Simulation-based calibration tests whether Bayesian credible ranks
are uniform when parameters and data are drawn from the stated prior and model.
Coverage experiments test frequentist intervals across controlled truth.
These checks address computation and calibration under the assumed model; they
do not prove that the assumed model describes nature.  Stress tests should
therefore inject departures that are scientifically plausible: non-Gaussian
line wings, spatially varying ratios, mismatched continua, correlated sky
residuals, or a contaminating transition at an uncertain redshift.

Blind analysis can reduce experimenter-induced tuning for delicate detections.
Negative controls, off-source apertures, inverted spectra, and line-free
wavelengths reveal false positives.  Split-data consistency is useful only when
the split isolates plausible failure modes, such as detector, epoch,
polarization, nod position, or array configuration.  Repeatedly choosing the
split that looks best creates another trials factor.

\begin{figure}[htbp]
\centering
\begin{tikzpicture}[font=\small,node distance=7mm and 9mm,
stage/.style={draw,rounded corners,fill=blue!4,align=center,minimum width=2.65cm,minimum height=8mm},
artifact/.style={draw,fill=green!5,align=center,minimum width=2.65cm,minimum height=8mm},
arr/.style={->,thick}]
\node[stage] (raw) {raw data +\\calibration context};
\node[stage,right=of raw] (reduce) {versioned\\reduction};
\node[artifact,right=of reduce] (spec) {spectrum/cube +\\mask + covariance};
\node[stage,below=of spec] (fit) {forward model +\\likelihood};
\node[artifact,left=of fit] (post) {posterior/likelihood +\\diagnostics};
\node[stage,left=of post] (pop) {selection-aware\\population model};
\node[artifact,below=of pop] (release) {tables + provenance +\\reproduction recipe};
\node[stage,below=of fit] (validate) {injection, recovery,\\predictive checks};
\draw[arr] (raw)--(reduce); \draw[arr] (reduce)--(spec); \draw[arr] (spec)--(fit);
\draw[arr] (fit)--(post); \draw[arr] (post)--(pop); \draw[arr] (pop)--(release);
\draw[arr] (fit)--(validate); \draw[arr] (validate)--(fit);
\draw[arr,dashed] (validate)--(release);
\end{tikzpicture}
\caption{A reproducible line-analysis graph.  Green artifacts are preserved
interfaces; blue stages are deterministic or explicitly stochastic workflows.
Validation is part of the inference loop, not a decorative final plot.}
\label{fig:reproducible-dag}
\end{figure}

\section{Provenance and machine-readable products}

Reproducibility requires more than a code repository.  Preserve immutable raw
identifiers and checksums; calibration reference names; software and dependency
versions; configuration files; random seeds; environment or container digest;
commands; and links between inputs and outputs.  Use stable object and line IDs,
machine-readable units, explicit missing-value semantics, and documented
coordinate and time frames.  FITS remains central in astronomy, while ECSV,
ASDF, HDF5, Parquet, and domain-specific models can preserve richer structure;
the correct choice is the one accompanied by a schema and validation.
The FAIR principles - findable, accessible, interoperable, and reusable - are a
helpful publication target, but ``accessible'' need not mean unrestricted when
privacy, proprietary periods, or facility policy apply \citep{Wilkinson2016}.
Metadata and a documented route to authorized access can remain public while
controlled data stay protected.

A line catalogue should minimally carry rest wavelength or frequency with source
and medium, transition identifier, observed centroid and frame, integrated flux,
signed estimator and uncertainty, width and its definition, continuum and
equivalent width, aperture, response or resolution, covariance reference,
detection procedure, quality mask, and provenance pointer.  Store posterior
samples or a reusable likelihood for nonlinear or multimodal results when
practical.  A symmetric error bar is not an adequate serialization of a bounded,
skewed, or mixture posterior.

Continuous integration should rebuild representative reductions or synthetic
tests, validate schemas and units, verify bibliography and documentation links,
and compare scientific outputs within justified tolerances.  Exact bitwise
reproduction is desirable for deterministic stages but can be brittle across
hardware; for stochastic inference, test seeds, distributional summaries, and
coverage.  Tag the release used for a publication and archive it with a durable
identifier.  A notebook can communicate exploration, but a parameterized,
non-interactive pipeline should produce the claimed tables and figures from
declared inputs.

A robust pipeline separates pure transformations from side effects.  A command
should accept declared inputs and configuration, write outputs to a new versioned
location, and emit a manifest containing checksums and environment identity.
Workflow engines can execute a directed acyclic graph, but the graph alone is
not provenance if calibration files are downloaded by mutable URL or notebooks
depend on hidden state.  Cache keys should include scientific configuration and
reference-data versions, not only source-code commits.  Manual quality decisions
belong in append-only, machine-readable annotations with reviewer, timestamp,
reason, and supersession history.

\section{A compact analysis contract}

Before publication, write a one-page contract stating: the primary observable;
the detector-space likelihood; response and background treatment; priors and
nuisance parameters; selection function; censoring policy; atomic/model versions;
uncertainty ledger; validation tests; and exact data/code release.  This contract
does not eliminate scientific judgment.  It makes judgment visible, testable,
and replaceable - the essential properties of a durable line measurement.
